% ECONOMICS_PROBLEM: {"id": "O33-298", "title": "AI and aggregate productivity", "jel_code": "O33", "source_id": "OP-0263"}
% !TeX program = xelatex
\documentclass[11pt,a4paper]{article}
\usepackage[margin=23mm]{geometry}
\usepackage{fontspec}
\setmainfont{Latin Modern Roman}
\usepackage{amsmath,amssymb,mathtools}
\usepackage{xcolor,enumitem,booktabs,longtable,array,xurl,hyperref}
\definecolor{muted}{HTML}{53636C}
\hypersetup{colorlinks=true,urlcolor=blue,linkcolor=blue}
\setlength{\parindent}{0pt}
\setlength{\parskip}{5pt}
\newcommand{\R}{\mathbb R}
\newcommand{\E}{\mathbb E}
\newcommand{\N}{\mathbb N}
\newcommand{\ind}{\mathbf 1}
\newcommand{\Var}{\operatorname{Var}}
\newcommand{\Cov}{\operatorname{Cov}}
\newcommand{\supp}{\operatorname{supp}}
\DeclareMathOperator*{\argmin}{arg\,min}
\DeclareMathOperator*{\argmax}{arg\,max}
\newcommand{\entrymeta}[3]{{\sffamily\small\color{muted}\textbf{\texttt{#1}}\quad|\quad #2\par #3\par}}
\newcommand{\Problemref}[1]{\texttt{#1}}
\newcommand{\JELref}[2]{\texttt{#2} (JEL \texttt{#1})}
\begin{document}
\section*{AI and aggregate productivity}
\textbf{Problem ID:} \texttt{O33-298}\quad \textbf{JEL:} \texttt{O33}\par
% BEGIN ECONOMICS STATEMENT
\label{problem:OP-0263}
{\sffamily\small\color{muted}Original subject group: Growth and productivity\par}
\label{definitions:problem:30007544}
\textbf{Audit classification:} Published research agenda. AI growth constraints are a published agenda; the defined capability and TFP contrast are editorial. Verification concerns the cited version, not first priority or proof that this exact editorial specification remains unsolved on October 8, 2026.
\subsection*{Definitions and precise research target}
\paragraph{Editorial mathematical specification.} Fix an economy with firms \(i\), occupations \(o\), and tasks \(j\). Let \(a_{ijt}\) denote use of a precisely defined AI capability, \(K^C_{it}\) complementary organizational capital, and \(K^A_{it}\) compute capital. Production combines task outputs through a declared aggregation technology; task demand, wages, entry, and investment respond in general equilibrium. Specify an AI-feasible counterfactual path \(a^1\) and a restricted path \(a^0\) in which that capability is unavailable, holding other primitive innovations fixed.
For measured aggregate total factor productivity \(A_t\) and horizon \(H\), define
\[
\Delta_H^{AI}=E[\log A_{t+H}(a^1)-\log A_{t+H}(a^0)].
\]
Identify a credible set for \(\Delta_H^{AI}\) using task experiments, adoption instruments, and independently measured complementary investment, with explicit assumptions linking microeconomic responses to equilibrium aggregation. Separate direct task efficiency from adoption selection, displaced labor, new task creation, and measurement of intangible capital. Impose observed constraints on compute, energy, managerial capacity, and diffusion rather than extrapolating laboratory task improvements mechanically. Validate transition predictions against subsequent adoption cohorts. The target concerns a declared capability, counterfactual, and horizon; forecasts about unspecified future intelligence or unconstrained autonomous research are not identified by current firm-level experiments.

An AI treatment must specify the capability, deployment permissions, price path, and adoption time. Define measured TFP by a fixed accounting convention, for example \(A_t=Y_t/(K_t^{\alpha}L_t^{1-\alpha})\), \(0<\alpha<1\), with positive measured inputs; specify whether compute and intangible capital enter \(K\). This residual need not equal a structural technology parameter when markups, utilization, and factor shares change. Both potential paths solve their respective resource-feasible equilibrium under a common law of non-AI innovations. The contrast is a general-equilibrium policy/capability effect, not a simple sum of experimental task gains. Micro-to-macro transport requires a stated task production function, substitutability, adoption response, and labor reallocation. Different AI definitions produce different estimands rather than interchangeable forecasts.
\subsection*{Variants and assumption changes}
The following are \emph{editorial variants}, not additional published open conjectures.
\begin{enumerate}
\item \emph{Fixed tasks and adoption.} Hold the task set and adoption shares fixed, identifying causal task-output gains from randomized access. Aggregate within a specified production function to obtain a conditional direct-productivity benchmark.
\item \emph{Endogenous task creation.} Permit new tasks, firm entry, and AI-assisted research to change technologies over time, while imposing measured energy and compute supply constraints. Estimate a finite-horizon equilibrium effect; current task experiments alone cannot identify future research acceleration.
\end{enumerate}
\subsection*{References and source audit}
\begin{itemize}
\item Erik Brynjolfsson; Anton Korinek; Ajay Agrawal. \emph{The Economics of Transformative AI: A Research Agenda}. 2024. Locator: Section3.1, Economic Growth, printed pp.9--10 (PDF pages9--10); October 23, 2024 draft. Primary text checked. \url{https://digitaleconomy.stanford.edu/app/uploads/2024/11/ETAI-White-Paper.pdf}.
\end{itemize}
AI growth constraints are a published agenda; the defined capability and TFP contrast are editorial.
\begin{enumerate}\raggedright
\item\hypertarget{definitions:source:30007544:1}{} Erik Brynjolfsson; Anton Korinek; Ajay Agrawal. 2024. \emph{The Economics of Transformative AI: A Research Agenda}. Section3.1, Economic Growth, printed pp.9--10 (PDF pages9--10); October23,2024 draft.
\url{https://digitaleconomy.stanford.edu/app/uploads/2024/11/ETAI-White-Paper.pdf}.
\end{enumerate}

\subsection*{Common conventions: Source mathematical conventions}

These conventions apply to each entry unless explicitly replaced.

\textbf{Models and identification.} A primitive model \(m\) specifies all probability laws, feasible choices, information, equilibrium and measurement rules used by its target. The admissible class \(\mathcal M\) is fixed independently of outcomes. If \(P_O(m)\) is its observable probability law and \(\tau(m)\) the target, its identified set at observable law \(P\) is
\[
 \mathcal I_\tau(P)=\{\tau(m):m\in\mathcal M,\ P_O(m)=P\}.
\]
Point identification means that this set is a singleton. An empty set signals incompatibility of data and assumptions. Coordinate bounds need not describe the joint set. Bounds are sharp only if every retained target is attainable or arbitrarily approachable by compatible models. Finite bounds require explicit restrictions on unobserved quantities; unrestricted confounding, hidden wealth, extrapolation or arbitrary equilibrium selection can yield uninformative sets.

\textbf{Causal interventions.} A potential outcome \(Y_i(p)\) is the outcome under a complete policy regime \(p\), including timing, financing, information and market spillovers. The target population and units are fixed before treatment. With interference, use the complete assignment vector or a justified exposure mapping; individual no-interference notation alone cannot identify an economy-wide intervention. A difference of mean potential outcomes does not require the cross-world joint distribution. Conditioning on post-treatment migration, survival, enrollment or employment changes the estimand and needs separate justification.

\textbf{Statistical statements.} An estimator, confidence set, forecast or test specifies a sampling law, model class and loss/coverage criterion. Uniform \((1-\alpha)\)-coverage means \(\inf_{m\in\mathcal M}\Pr_m\{\tau(m)\in C_n\}\geq1-\alpha\), or an explicitly asymptotic version. The whole parameter space trivially has coverage, so informativeness requires a stated width, risk or power objective. A forecasting improvement is relative to a named benchmark, information set and evaluation population.

\textbf{Equilibrium and optimization.} An equilibrium comprises feasible plans, optimality, beliefs where needed, budget constraints and all market/resource clearing. The primitives or a stated benchmark define each correspondence \(\mathcal E(p;m)\). Nonemptiness is assumed or proved, never inferred just from compact policy sets. A selected-equilibrium target uses a declared selection rule; a robust target quantifies over all permitted equilibria. Use suprema or infima unless attainment is established. An \(\varepsilon\)-optimal policy is feasible and within \(\varepsilon\) of the finite optimum value. Report an empty feasible set separately.

\textbf{Welfare and accounting.} Normative utility, interpersonal normalization, social weights, numeraire, discounting and the use of public receipts are primitives. Behavioral utility need not coincide with the welfare criterion. Household welfare includes ownership income and rents once; adding profits again to compensating variation generally double-counts them. A monetary equivalent is defined by an explicit transfer and price convention, with monotonicity and range conditions. Average effects, mechanism contributions, transfers and welfare losses are different targets. Infinite sums require convergence and dynamic feasible plans require terminal or no-Ponzi conditions.

\textbf{Source versions.} A working-paper date, revision date and journal publication year are distinguished. A DOI for a journal version is not automatically the DOI of an earlier manuscript. Locators refer to the stated version and printed pages where specified. A metadata-only check does not verify a claim about a full-text conclusion. Every entry's source subsection narrows the provenance claim accordingly.
% END ECONOMICS STATEMENT
\end{document}
